\subsection{Arsitektur Sistem}\label{rancangan-solusi-arsitektur-sistem}

Gambaran high-level:
\begin{itemize}
	\item Patient device (Android)
	\item Optional: wearable $\rightarrow$ Android
	\item Local DB + batching engine
	\item Encryption module
	\item IPFS cluster
	\item IOTA node + smart contract DecMed (extended)
	\item Existing DecMed components (proxy, CapBAC, PRE) di mana PGHD diintegrasikan
\end{itemize}



Interval pengambilan data dilakukan setiap interval waktu pada tabel~\ref{health-connect-get-data-interval} \parencite{android_hc_write}.

\begin{table}[htbp]
	\centering
	\caption{Panduan Interval Pengambilan Data}\label{health-connect-get-data-interval}
	\renewcommand{\arraystretch}{1.25}
	\setlength{\tabcolsep}{8pt}
	\rowcolors{2}{white}{white}
	\begin{tabularx}{\textwidth}{>{\raggedright\arraybackslash}p{0.30\textwidth}>{\raggedright\arraybackslash}X}
		\rowcolor[HTML]{DDEBF7}
		\textbf{Tipe Data} & \textbf{Ekspektasi (Frekuensi Penulisan)} \\
		\midrule
		Steps & Setiap 1 menit \\
		\midrule
		StepsCadence & Setiap 1 menit \\
		\midrule
		WheelchairPushes & Setiap 1 menit \\
		\midrule
		ActiveCaloriesBurned & Setiap 15 menit \\
		\midrule
		TotalCaloriesBurned & Setiap 15 menit \\
		\midrule
		Distance & Setiap 1 menit \\
		\midrule
		ElevationGained & Setiap 1 menit \\
		\midrule
		FloorsClimbed & Setiap 1 menit \\
		\midrule
		HeartRate & 4 kali per menit \\
		\midrule
		HeartRateVariabilityRmssd & Setiap 1 menit \\
		\midrule
		RespiratoryRate & Setiap 1 menit \\
		\midrule
		OxygenSaturation & Setiap jam \\
		\bottomrule
	\end{tabularx}
\end{table}



Interval 15 menit dipilih merujuk pada panduan penulisan data Health Connect API yang merekomendasikan interval tersebut untuk tipe data fisiologis
Durasi yang digunakan untuk mengirimkan data dari klien ke sistem DecMed adalah setiap 15 menit sekali seperti rekomendasi Google HealthConnect API \parencite{android_hc_write}.